\section{Security Team Maturity：从救火队到嵌入式治理}

团队成熟度不是 headcount 曲线。一个 5 人团队可能拥有清晰 authority 和高质量自动化；一个 100 人团队也可能只是接收告警、转发 ticket。成熟的 security organization 会逐步把 threat signal、incident workflow、risk acceptance 和 product decision 嵌入业务系统，使安全不再是发布前的临时审批。

\subsection{四个阶段与错误捷径}

Reactive 阶段依赖少数专家救火；program 阶段建立 intake、on-call、policy 和 repeatable controls；platform 阶段把 identity、telemetry、automation 和 guardrails 做成共享能力；embedded 阶段让 product、finance、legal、HR 和 communications 在日常决策中承担自己的安全责任。阶段不会自动前进，尤其不能用购买产品假装完成组织变革。

\lecturefigure{13-security-maturity.png}{安全团队从 reactive response 走向 embedded governance。}{本地字幕 00:00--03:20；31:20--33:40；概念重绘。}

\begin{knowledgebox}{读图：成熟度的输出不是更多控制，而是更短反馈}
Reactive 只能在事件后响应；program 建立角色和运行节奏；platform 让控制可复用、可观察；embedded 把产品 ownership、risk acceptance 与持续改进送回业务。判断成熟度时先看 handoff 是否明确、决策是否留痕、例外是否过期、postmortem 是否改变系统，而不是看 policy 数量或采购金额。
\end{knowledgebox}

\begin{warningbox}{“Security owns risk” 是危险表述}
安全团队通常负责识别、分析、建议和监控，但业务负责人拥有产品收益与残余风险，法务负责法律判断，高管和董事会承担重大企业决策。把所有风险都交给 CSO，会让其他负责人在平时缺席、在事故后寻找单点责任人。
\end{warningbox}

\begin{importantbox}{成熟团队的最小接口}
至少定义：统一 intake、severity taxonomy、incident commander、legal/comms escalation、evidence repository、materiality workstream、exception register、executive exercise 和 post-incident owner。每个接口都应有响应时间、替补角色与可审计输出。
\end{importantbox}

\subsection{本章小结}

团队成长的目标是把个人经验编码为组织接口。越接近 embedded 阶段，越少出现“安全团队知道但高管不知道”“法务判断了但没有技术事实”或“事件结束却无人修改系统”的断裂。

